% Source PDF: source/普通高考/2006/2006全国2文(黑龙江,吉林,贵州,新疆,内蒙古,青海,云南,西藏,宁夏,,甘肃).pdf
% Page: 1
% PDF embedded bitmap attempt: no matching image objects; the page is vector text/line art.
% Grid evidence: tmp/repaint_grid_2006_n2l/page1_grid100.png and tmp/repaint_grid_2006_n2l/q16_block_grid50.png
% Page crop evidence: tmp/repaint_crops_2006_n2l/q16_orig.png, cropped from the no-grid page render.
% Elements: x-axis with origin 0 and a break before 1000; y-axis; six histogram bars over [1000,1500), [1500,2000), [2000,2500), [2500,3000), [3000,3500), [3500,4000); dashed horizontal guides at 0.0001 through 0.0005; labels 频率/组距 and 月收入（元）.
% Relations: bar heights are 0.0002, 0.0004, 0.0005, 0.0005, 0.0003, 0.0001 respectively, read from the source figure and consistent with the stated frequency-density histogram.
% solid segments: bar outlines, x/y axis arrows, tick marks, and the x-axis break stroke.
% dashed segments: horizontal reference lines at y=0.0001, 0.0002, 0.0003, 0.0004, 0.0005.
% labels: y label sits above-left of the vertical axis; x label sits beyond the 4000 tick; tick labels stay below/left of axes and clear the bars.
% x display map: 0 -> 0, then each true group boundary is shifted left by 500
% so the displayed 0-to-1000 break span is exactly one 500-unit group width.
\begin{tikzpicture}[baseline]
\begin{axis}[exam statistical histogram,
    width=12.5cm,
    axis lines=none,
    clip=false,
    xmin=-180,
    xmax=4000,
    ymin=-0.00009,
    ymax=0.00059,
    ticks=none,
    x tick label as interval=false,
]
    \draw[dashed,line width=0.55pt] (axis cs:0,0.0001) -- (axis cs:3500,0.0001);
    \draw[dashed,line width=0.55pt] (axis cs:0,0.0002) -- (axis cs:3500,0.0002);
    \draw[dashed,line width=0.55pt] (axis cs:0,0.0003) -- (axis cs:3500,0.0003);
    \draw[dashed,line width=0.55pt] (axis cs:0,0.0004) -- (axis cs:3500,0.0004);
    \draw[dashed,line width=0.55pt] (axis cs:0,0.0005) -- (axis cs:3500,0.0005);

    \addplot[
        ybar interval,
        mark=none,
        draw=black,
        fill=white,
        line width=0.7pt,
    ] coordinates {
        (500,0.0002)
        (1000,0.0004)
        (1500,0.0005)
        (2000,0.0005)
        (2500,0.0003)
        (3000,0.0001)
        (3500,0)
    };

    \draw[->,line width=0.75pt] (axis cs:0,0) -- (axis cs:3800,0);
    \draw[->,line width=0.75pt] (axis cs:0,0) -- (axis cs:0,0.00058);

    % Keep the two-tooth break broad across the compressed first group, but
    % shallow enough that it stays visually attached to the baseline.
    \draw[white,line width=1.8pt] (axis cs:60,0) -- (axis cs:440,0);
    \examstatxbreak[line width=0.75pt]{250}{0};

    \examstatylabel{\(\dfrac{\text{频率}}{\text{组距}}\)};
    % Place the x-axis title on its own baseline row so the complete tick-label
    % bbox remains measurable and the title stays inside the figure boundary.
    \examstatxlabel{月收入（元）};

    \draw[line width=0.45pt] (axis cs:0,0) -- (axis cs:0,-0.000014);
    \node[anchor=east,inner sep=0pt,font=\normalsize,xshift=-4pt] at (axis cs:0.000000,0) {0};
    \draw[line width=0.45pt] (axis cs:500,0) -- (axis cs:500,-0.000014);
    \node[anchor=north,inner sep=0pt,font=\normalsize,yshift=-2pt] at (axis cs:500,-0.000014) {\(1000\)};
    \draw[line width=0.45pt] (axis cs:1000,0) -- (axis cs:1000,-0.000014);
    \node[anchor=north,inner sep=0pt,font=\normalsize,yshift=-2pt] at (axis cs:1000,-0.000014) {\(1500\)};
    \draw[line width=0.45pt] (axis cs:1500,0) -- (axis cs:1500,-0.000014);
    \node[anchor=north,inner sep=0pt,font=\normalsize,yshift=-2pt] at (axis cs:1500,-0.000014) {\(2000\)};
    \draw[line width=0.45pt] (axis cs:2000,0) -- (axis cs:2000,-0.000014);
    \node[anchor=north,inner sep=0pt,font=\normalsize,yshift=-2pt] at (axis cs:2000,-0.000014) {\(2500\)};
    \draw[line width=0.45pt] (axis cs:2500,0) -- (axis cs:2500,-0.000014);
    \node[anchor=north,inner sep=0pt,font=\normalsize,yshift=-2pt] at (axis cs:2500,-0.000014) {\(3000\)};
    \draw[line width=0.45pt] (axis cs:3000,0) -- (axis cs:3000,-0.000014);
    \node[anchor=north,inner sep=0pt,font=\normalsize,yshift=-2pt] at (axis cs:3000,-0.000014) {\(3500\)};
    \draw[line width=0.45pt] (axis cs:3500,0) -- (axis cs:3500,-0.000014);
    \node[anchor=north,inner sep=0pt,font=\normalsize,yshift=-2pt] at (axis cs:3500,-0.000014) {\(4000\)};

    \draw[line width=0.45pt] (axis cs:0,0.0001) -- (axis cs:-55,0.0001);
    \node[anchor=east,inner sep=0pt,font=\normalsize,xshift=-4pt] at (axis cs:-55.000000,0.0001) {0.0001};
    \draw[line width=0.45pt] (axis cs:0,0.0002) -- (axis cs:-55,0.0002);
    \node[anchor=east,inner sep=0pt,font=\normalsize,xshift=-4pt] at (axis cs:-55.000000,0.0002) {0.0002};
    \draw[line width=0.45pt] (axis cs:0,0.0003) -- (axis cs:-55,0.0003);
    \node[anchor=east,inner sep=0pt,font=\normalsize,xshift=-4pt] at (axis cs:-55.000000,0.0003) {0.0003};
    \draw[line width=0.45pt] (axis cs:0,0.0004) -- (axis cs:-55,0.0004);
    \node[anchor=east,inner sep=0pt,font=\normalsize,xshift=-4pt] at (axis cs:-55.000000,0.0004) {0.0004};
    \draw[line width=0.45pt] (axis cs:0,0.0005) -- (axis cs:-55,0.0005);
    \node[anchor=east,inner sep=0pt,font=\normalsize,xshift=-4pt] at (axis cs:-55.000000,0.0005) {0.0005};
\end{axis}
\end{tikzpicture}
